\documentclass[12pt]{article}
\usepackage[utf8]{inputenc}
\usepackage[T1]{fontenc}
\usepackage{hyperref}
\title{Birrpayi Language Resources for K-12 Curriculum Development}
\author{Compiled by Claude Sonnet 5.5 (claude-sonnet-5-5) for Andr\'as Kornai}
\date{}
\begin{document}
\maketitle

\section{Introduction}
Birrpayi (ISO 639-3 xbj; spellings include Birrbay and Birpai) is an extinct Australian language of the New South Wales mid-north coast. Glottolog lists it as extinct and notes, citing Lissarrague (2010), that it was close enough to count as a dialect of Gathang (Worimi, kda). Wafer and Lissarrague (2008) treat Gadhang, Warrimay, Guringay and Birrbay as dialects of one ``Lower North Coast Language''. W. J. Enright found four elderly speakers of Gathang at Wauchope in 1932. It should not be confused with the Queensland languages Biri or Pirriya. The revival work that exists concerns Gathang. I found no Birrpayi-specific school or revival materials.

\section{Primary Educational Resources}
None found for Birrpayi specifically.

\section{Linguistic and Reference Materials}
\begin{itemize}
\item Glottolog record for Birrpayi (xbj). \url{https://glottolog.org/resource/languoid/iso/xbj}
\item Wikipedia, ``Gathang language'', which covers the Birrbay, Guringay and Warrimay dialects. \url{https://en.wikipedia.org/wiki/Gathang_language}
\item Wikipedia, ``Birrbay''. \url{https://en.wikipedia.org/wiki/Birrbay}
\item AIATSIS AUSTLANG entry E3. \url{https://collection.aiatsis.gov.au/austlang/language/e3}
\end{itemize}

\section{Religious and Community Resources}
None found for Birrpayi specifically. Birrbay (Birpai) community research material is held at the Mid North Coast library research collection. \url{https://mnclibrary.org.au/research/birpai-studies/}

\section{Digital and Audio Resources}
None found.

\section{Curriculum Development Context}
Because Birrpayi is so close to Gathang, a heritage curriculum would probably draw on the Gathang revival materials, which are not documented here. Any such curriculum should be developed with Birrbay and Worimi community organisations. Little Birrpayi-specific material exists in what I found.
\end{document}